% RESEARCH_DRAFT_SHA256: cbb7f1af1e24d63bfa83f76b3c2a91766b0298d42604b8eeb2f1142d5e96e213
% Layout and source-reference wrappers expanded for the web renderer.
\documentclass[11pt]{article}
\usepackage[margin=25mm]{geometry}
\usepackage{fontspec,unicode-math}
\setmainfont{Latin Modern Roman}
\setmathfont{Latin Modern Math}
\usepackage{amsmath,amssymb,mathtools,xurl,hyperref}
\hypersetup{colorlinks=true,urlcolor=blue,linkcolor=blue}
\setcounter{secnumdepth}{0}
\setlength{\parindent}{0pt}
\setlength{\parskip}{5pt}
\setlength{\emergencystretch}{3em}
\newcommand{\sref}[2]{\hyperlink{source:#1:#2}{[S#2]}}
\newcommand{\cataloguescope}[1]{\paragraph{Source record.}#1}
\newcommand{\floor}[1]{\left\lfloor#1\right\rfloor}
\begin{document}
\section*{Research-attempt provenance}
This standalone research attempt retains the complete original definition body below,
including its historical source and preparation notes. Those notes describe the
earlier catalogue preparation; the new attempt and its verification are identified
separately after that body. No original ID, catalogue status, or source statement is changed.

% UnsolvedMath ID 20002181; source catalogue code AIM-INFRASTRUCTURE-0081
\section{Certifying offline PreTeXt slides on removable media}\label{20002181}
{\small UnsolvedMath ID 20002181 · Computer Science}
\subsection{Definitions and mathematical statement}

\paragraph{Shared notation.}Standard number systems are $\mathbb N,\mathbb Z,\mathbb Q,\mathbb R,\mathbb C$. Specialized notation and index ranges follow the cited source and the definitions below.


An offline slide presentation must make its required content and assets available without network access. A removable drive is a transport medium and does not by itself supply a rendering environment. Portability depends on formats, dependencies, and the software available at the destination.
\paragraph{Record scope and qualifications.}
The catalogue title is retained as metadata. Named certificates, contracts, bounds or frameworks in that title are not asserted as results of the original source extract.
The registry marks the fragment reconstructed; no required format, offline guarantee, or target device is specified.

Catalogue screening limitation: PreTeXt software design, documentation, or publication feature request; no mathematical problem is specified by the supplied passage.

The source snapshot labels this record partially solved; this is catalogue metadata rather than independent verification of a mathematical result or the current completion of a project.
The source text below is rendered literally where necessary; damaged formulas and composite extracts are not silently repaired into a new mathematical claim.
The supplied catalogue marks this formulation as an unverified reconstruction; missing assumptions are not certified by the source link.
\paragraph{Statement recorded in the supplied source.}
Other use cases: slides on a thumb drive \href{https://github.com/PreTeXtBook/community/wiki/Moderated-Problem-Session}{[S2]}
\paragraph{Source status.}
The supplied catalogue labels this source record partially solved; that metadata was checked against the supplied snapshot on 2026-10-11. The label is preserved for provenance and is not an independent verification of present open status or mathematical correctness.
\subsection{Short English statement}
Consider portable, offline slide presentations as another publication use case.
\subsection{Sources}
\begin{itemize}
\item[S1] Ulam AI and the UnsolvedMath Contributors, UnsolvedMath: A Curated Collection of Open Mathematics Problems, record 20002181 (AIM-INFRASTRUCTURE-0081). Catalogue curation and metadata: CC BY 4.0. The source extract is retained; definitions, English summaries and qualifications are editorial additions.
\url{https://huggingface.co/datasets/ulamai/UnsolvedMath}.
\item[S2] AIM Workshop Problem Lists, source linked by catalogue record AIM-INFRASTRUCTURE-0081.
\url{https://github.com/PreTeXtBook/community/wiki/Moderated-Problem-Session}.
\end{itemize}
\paragraph{Preparation and review.}
The original source and UnsolvedMath contributors retain attribution for the source extract. Codex (GPT-6; exact serving revision not exposed) assisted with definitions, English summaries, scope notes and formatting. No new human mathematical authorship or independently verified solution is asserted. Permanent source ID 20002181 and code AIM-INFRASTRUCTURE-0081 checked against the complete supplied registry and current repository. Distinct upstream records, including repeated questions and broad or damaged extracts, are retained for complete coverage of the source CS category. A repeated record is not a new independent conjecture. Complete mathematical-equivalence and current-literature review has not been performed. Preparation and final review follow contribution prompts C and D, with source gaps disclosed. Structural and rendering checks do not establish a new theorem or open status.
\label{end:20002181}

\section{Literature and context check}
The primary community-session source was consulted. The record is a broad software-design request, and no mathematical theorem follows from its editorial catalogue title. The finite model below is explicitly auxiliary; it does not certify an actual implementation or current feature status.
\section{Research attempt: An offline slide test must cover reachable requests}
A single successful slide-show trace is insufficient evidence of offline completeness. Consider a finite state machine with initial slide $s$, a next action to a fully local slide $t$, and a notes action to state $u$ that requests a missing remote resource. Testing only next succeeds, while notes fails offline. Both actions start from the same initial slide, so this is not a difference of document versions.

For a specified finite deterministic runtime model, annotate each transition with the resources it requests and let $L$ be the locally supplied resource set. Explore all states reachable from the initial state. The offline-request property holds exactly when every transition out of every reachable state requests only elements of $L$. Necessity: a reachable bad edge has a path to its source and then an action producing a missing request. Sufficiency: by induction over an arbitrary action sequence, each traversed source state is reachable and each requested resource belongs to $L$. A breadth-first search can retain a shortest predecessor path, giving a concrete failing action sequence when it finds a bad edge. Runtime capability failures must be modeled as bad transitions too; merely finding a file is not enough.

The finite-state and complete-request assumptions are material. If a program computes a resource name from unrestricted input, testing finitely many traces cannot enumerate all possibilities without further analysis. A model with a finite declared input alphabet and state set can be checked exhaustively; the real browser, removable-media paths, permissions, and script behavior were not modeled or tested here. The argument is therefore a reproducible criterion for a bounded profile, not certification of actual slides on a thumb drive.
\section{Verification and adversarial audit}
The construction and counterexample are checked in the written finite model. No actual PreTeXt build, browser, EPUB reader, or publishing service was tested. All runtime and identity assumptions used by the proof are stated explicitly.
\paragraph{Audit corrections.} One observed successful path is not extrapolated to all interactions. The positive proof requires a complete finite transition model and separately represented runtime failures.
\paragraph{External computation.} None; symbolic arguments and boundary checks only.
\section{Scope and remaining gaps}
This is a restricted route investigation. The full feature requires implementation, supported-version decisions, conformance tests, and user-facing evaluation that were not performed here. The original agenda remains unresolved, and no novelty or completed software feature is claimed.
\paragraph{Integration review.} Codex (GPT-6; exact serving revision not exposed) reviewed the written arguments and expanded layout and source-reference wrappers. Repository metadata and the production web renderer were checked using Python 3.13 and JavaScript with MathJax 3.2.2. This model-assisted review is not independent human proof verification; no PDF compilation or formal proof checking is asserted.

\section{Completion Estimate}
\noindent\textbf{COMPLETION: 0\%}
The calculations below the preserved source statement investigate a concrete route,
but establish no substantiated new progress toward the entire original target.
Known special cases and failed sufficient conditions are not credited as a solution.
\section{References}
The original source references and catalogue attribution are preserved in the
Sources section above. The following URLs were actually consulted for this attempt;
retention of an additional original reference does not claim fresh consultation of it.
\begin{itemize}
\item \url{https://github.com/PreTeXtBook/community/wiki/Moderated-Problem-Session} (consulted 2026-10-11). 
\end{itemize}
\section{Model and contributor disclosure}
OpenAI gpt-6-astra, requested reasoning configuration ultra, generated this research
attempt on 2026-10-11. The exact serving revision is not exposed. Codex GPT-6
performed orchestration and repository-format review; its exact serving revision
is also not exposed. The mathematical route, boundary checks and final adversarial
self-audit followed contribution prompts A then B in
\texttt{docs/contribution\_prompts.md}, after reading \texttt{CONTRIBUTING.md}.
Python 3.14.2 assembled the standalone documents; any mathematical computation is
specified in Verification. No theorem prover was used. No human mathematical
authorship, substantial human proof correction, or independent human proof checking
is asserted. Model-assisted review, source retrieval, and structural checks do not
establish correctness or novelty. Full private reasoning transcripts are unavailable
as reproducibility artifacts; the complete arguments, inputs and outcomes actually
used are stated here or in the companion batch audit.
\section{Final status}
Unresolved as a new research contribution. This is a concrete route investigation,
not a claim of new complete resolution or a current-status certification. A known
answer reproduced for context does not become a new contribution. The catalogue's
source-status label is preserved; this attempt status does not assert that a
source-labelled solved question is open.
\end{document}
